\ifdefined\gkver\else\def\gkver{student}\fi
\documentclass[\gkver]{gaokaozhenti}
\begin{document}

\chapter{2021年河北}

\begin{enumerate}
\item
%% number: 1
%% paperName: 2021年河北高考第 1 题 4 分
%% typeId: 1
%% type: 单选题
%% chapter: 原子和原子核
%% point: 半衰期
%% method: 无
%% score: 4
%% degree: 750
%% duplicateId: 0
%% body:
银河系中存在大量的铝同位素 \ce{^{26}Al}，\ce{^{26}Al} 核 $\beta$ 衰变的衰变方程为 $\ce{^{26}_{13}Al -> ^{26}_{12}Mg + ^{0}_{1}e}$，测得 \ce{^{26}Al} 核的半衰期为 72 万年，下列说法正确的是\xzanswer{C}
\fourchoices[answer=C]
{\ce{^{26}Al} 核的质量等于 \ce{^{26}Mg} 核的质量}
{\ce{^{26}Al} 核的中子数大于 \ce{^{26}Mg} 核的中子数}
{将铝同位素 \ce{^{26}Al} 放置在低温低压的环境中，其半衰期不变}
{银河系中现有的铝同位素 \ce{^{26}Mg} 将在 144 万年后全部衰变为 \ce{^{26}Mg}}

%% answer: C
%% memo:
\memoanswer{
A．\ce{^{26}_{13}Al} 和 \ce{^{26}_{12}Mg} 的质量数均为 26 相等，但是二者原子核中的质子数和中子数不同，所以质量不同，A 错误；

B．\ce{^{26}_{13}Al} 核的中子数为 $16-13=13$ 个，\ce{^{26}_{12}Mg} 核的中子数为 $26-12=14$ 个，B 错误；

C．半衰期是原子核固有的属性，与外界条件无关，C 正确；

D．质量为 $m$ 的 \ce{^{26}_{13}Al} 的半衰期为 72 万年，经过 $144=2\times72$ 万年为 2 个半衰期，剩余质量为 $\frac{1}{4}m$，不会全部衰变为 \ce{^{26}_{12}Mg}，D 错误。

故选 C。
}

\item
%% number: 2
%% paperName: 2021年河北高考第 2 题 4 分
%% typeId: 1
%% type: 单选题
%% chapter: 曲线运动
%% point: 平抛运动
%% method: 无
%% score: 4
%% degree: 750
%% duplicateId: 0
%% body:
铯原子钟是精确的计时仪器，图~\ref{2021河北02a}中铯原子从 O 点以 $100\Ums$ 的初速度在真空中做平抛运动，到达竖直平面 MN 所用时间为 $t_{1}$；图~\ref{2021河北02b}中铯原子在真空中从 P 点做竖直上抛运动，到达最高点 Q 再返回 P 点，整个过程所用时间为 $t_{2}$，O 点到竖直平面 MN、P 点到 Q 点的距离均为 $0.2\Um$，重力加速度取 $g=10\Umsq$，则 $t_{1}:t_{2}$ 为\xzanswer{C}
\twopicture[labelA=2021河北02a, labelB=2021河北02b, numA=1, numB=2, showlabel=false, w=4.2cm]{figs/02a.png}{figs/02b.png}
\fourchoices[answer=C]
{$100:1$}
{$1:100$}
{$1:200$}
{$200:1$}

%% answer: C
%% memo:
\memoanswer{
铯原子做平抛运动，水平方向上做匀速直线运动，即
$x=vt_{1}$
解得
$t_{1}=\frac{x}{v}=\frac{0.2}{100}\Us$
铯原子做竖直上抛运动，抛至最高点用时 $\frac{t_{2}}{2}$，逆过程可视为自由落体，即
$x=\frac{1}{2}g\left(\frac{t_{2}}{2}\right)^{2}$
解得
$t_{2}=\sqrt{\frac{8x}{g}}=\sqrt{\frac{8\times0.2}{10}}=0.4\Us$
则
$\frac{t_{1}}{t_{2}}=\frac{\frac{0.2}{100}}{0.4}=\frac{1}{200}$
故选 C。
}

\item
%% number: 3
%% paperName: 2021年河北高考第 3 题 4 分
%% typeId: 1
%% type: 单选题
%% chapter: 运动和力
%% point: 力学单位制
%% method: 无
%% score: 4
%% degree: 650
%% duplicateId: 0
%% body:
普朗克常量 $h=6.626\times10^{-34}\UJcdots$，光速为 $c$，电子质量为 $m_{e}$，则 $\frac{h}{m_{e}c}$ 在国际单位制下的单位是\xzanswer{B}
\fourchoices[answer=B]
{$\UJs$}
{$\Um$}
{$\mathrm{J}\cdot\mathrm{m}$}
{$\Ums$}

%% answer: B
%% memo:
\memoanswer{
根据 $\frac{h}{m_{e}c}$ 可得它们的单位为：
$\frac{\mathrm{J}\cdot\mathrm{s}}{\mathrm{kg}\cdot\mathrm{m}/\mathrm{s}}=\frac{\mathrm{N}\cdot\mathrm{m}\cdot\mathrm{s}}{\mathrm{kg}\cdot\mathrm{m}/\mathrm{s}}=\frac{\mathrm{kg}\cdot\mathrm{m}/\mathrm{s}^{2}\cdot\mathrm{m}\cdot\mathrm{s}}{\mathrm{kg}\cdot\mathrm{m}/\mathrm{s}}=\mathrm{m}$
故选 B。
}

\item
%% number: 4
%% paperName: 2021年河北高考第 4 题 4 分
%% typeId: 1
%% type: 单选题
%% chapter: 曲线运动
%% point: 人造地球卫星
%% method: 无
%% score: 4
%% degree: 550
%% duplicateId: 0
%% body:
“祝融号”火星车登陆火星之前，“天问一号”探测器沿椭圆形的停泊轨道绕火星飞行，其周期为 2 个火星日，假设某飞船沿圆轨道绕火星飞行，其周期也为 2 个火星日，已知一个火星日的时长约为一个地球日，火星质量约为地球质量的 0.1 倍，则该飞船的轨道半径与地球同步卫星的轨道半径的比值约为\xzanswer{D}
\fourchoices[answer=D]
{$\sqrt[3]{4}$}
{$\sqrt[3]{\frac{1}{4}}$}
{$\sqrt[3]{\frac{5}{2}}$}
{$\sqrt[3]{\frac{2}{5}}$}

%% answer: D
%% memo:
\memoanswer{
绕中心天体做圆周运动，根据万有引力提供向心力，可得
$\frac{GMm}{R^{2}}=m\frac{4\pi^{2}}{T^{2}}R$
则
$T=\sqrt{\frac{4\pi^{2}R^{3}}{GM}}$，$R=\sqrt[3]{\frac{GMT^{2}}{4\pi^{2}}}$
由于一个火星日的时长约为一个地球日，火星质量约为地球质量的 0.1 倍，则飞船的轨道半径
$R_{\text{飞}}=\sqrt[3]{\frac{GM_{\text{火}}(2T)^{2}}{4\pi^{2}}}=\sqrt[3]{\frac{G\times0.1M_{\text{地}}\times4\times\frac{4\pi^{2}R_{\text{同}}^{3}}{GM_{\text{地}}}}{4\pi^{2}}}=\sqrt[3]{\frac{2}{5}}R_{\text{同}}$
则
$\frac{R_{\text{飞}}}{R_{\text{同}}}=\sqrt[3]{\frac{2}{5}}$
故选 D。
}

\item
%% number: 5
%% paperName: 2021年河北高考第 5 题 4 分
%% typeId: 1
%% type: 单选题
%% chapter: 磁场
%% point: 洛伦兹力
%% method: 无
%% score: 4
%% degree: 450
%% duplicateId: 0
%% body:
如图，距离为 $d$ 的两平行金属板 P、Q 之间有一匀强磁场，磁感应强度大小为 $B_{1}$，一束速度大小为 $v$ 的等离子体垂直于磁场喷入板间，相距为 $L$ 的两光滑平行金属导轨固定在与导轨平面垂直的匀强磁场中，磁感应强度大小为 $B_{2}$，导轨平面与水平面夹角为 $\theta$，两导轨分别与 P、Q 相连，质量为 $m$、电阻为 $R$ 的金属棒 ab 垂直导轨放置，恰好静止，重力加速度为 $g$，不计导轨电阻、板间电阻和等离子体中的粒子重力，下列说法正确的是\xzanswer{B}
\onepicture[w=7.0cm]{figs/05.png}
\fourchoices[answer=B]
{导轨处磁场的方向垂直导轨平面向上，$v=\frac{mgR\sin\theta}{B_{1}B_{2}Ld}$}
{导轨处磁场的方向垂直导轨平面向下，$v=\frac{mgR\sin\theta}{B_{1}B_{2}Ld}$}
{导轨处磁场的方向垂直导轨平面向上，$v=\frac{mgR\tan\theta}{B_{1}B_{2}Ld}$}
{导轨处磁场的方向垂直导轨平面向下，$v=\frac{mgR\tan\theta}{B_{1}B_{2}Ld}$}

%% answer: B
%% memo:
\memoanswer{
等离子体垂直于磁场喷入板间时，根据左手定则可得金属板 Q 带正电荷，金属板 P 带负电荷，则电流方向由金属棒 a 端流向 b 端。等离子体穿过金属板 P、Q 时产生的电动势 $U$ 满足
$q\frac{U}{d}=qB_{1}v$
由欧姆定律 $I=\frac{U}{R}$ 和安培力公式 $F=BIL$ 可得
$F_{\text{安}}=B_{2}L\times\frac{U}{R}=\frac{B_{2}B_{1}Lvd}{R}$
再根据金属棒 ab 垂直导轨放置，恰好静止，可得
$F_{\text{安}}=mg\sin\theta$
则
$v=\frac{mgR\sin\theta}{B_{1}B_{2}Ld}$
金属棒 ab 受到的安培力方向沿斜面向上，由左手定则可判定导轨处磁场的方向垂直导轨平面向下。故选 B。
}

\item
%% number: 6
%% paperName: 2021年河北高考第 6 题 4 分
%% typeId: 1
%% type: 单选题
%% chapter: 功和能
%% point: 机械能守恒定律
%% method: 无
%% score: 4
%% degree: 550
%% duplicateId: 0
%% body:
一半径为 $R$ 的圆柱体水平固定，横截面如图所示，长度为 $\pi R$、不可伸长的轻细绳，一端固定在圆柱体最高点 P 处，另一端系一个小球，小球位于 P 点右侧同一水平高度的 Q 点时，绳刚好拉直，将小球从 Q 点由静止释放，当与圆柱体未接触部分的细绳竖直时，小球的速度大小为（重力加速度为 $g$，不计空气阻力）\xzanswer{A}
\onepicture[w=6.0cm]{figs/06.png}
\fourchoices[answer=A]
{$\sqrt{(2+\pi)gR}$}
{$\sqrt{2\pi gR}$}
{$\sqrt{2(1+\pi)gR}$}
{$2\sqrt{gR}$}

%% answer: A
%% memo:
\memoanswer{
小球下落的高度为
$h=\pi R-\frac{\pi}{2}R+R=\frac{\pi+2}{2}R$
小球下落过程中，根据动能定理有
$mgh=\frac{1}{2}mv^{2}$
综上有
$v=\sqrt{(2+\pi)gR}$
故选 A。
}

\item
%% number: 7
%% paperName: 2021年河北高考第 7 题 4 分
%% typeId: 1
%% type: 单选题
%% chapter: 电磁感应
%% point: 切割磁感线电动势
%% method: 无
%% score: 4
%% degree: 450
%% duplicateId: 0
%% body:
如图，两光滑导轨水平放置在竖直向下的匀强磁场中，磁感应强度大小为 $B$，导轨间距最窄处为一狭缝，取狭缝所在处 O 点为坐标原点，狭缝右侧两导轨与 $x$ 轴夹角均为 $\theta$，一电容为 $C$ 的电容器与导轨左端相连，导轨上的金属棒与 $x$ 轴垂直，在外力 $F$ 作用下从 O 点开始以速度 $v$ 向右匀速运动，忽略所有电阻，下列说法正确的是\xzanswer{A}
\onepicture[w=6.5cm]{figs/07.png}
\fourchoices[answer=A]
{通过金属棒的电流为 $2BCv^{2}\tan\theta$}
{金属棒到达 $x_{0}$ 时，电容器极板上的电荷量为 $BCvx_{0}\tan\theta$}
{金属棒运动过程中，电容器的上极板带负电}
{金属棒运动过程中，外力 $F$ 做功的功率恒定}

%% answer: A
%% memo:
\memoanswer{
C．根据楞次定律可知电容器的上极板应带正电，C 错误；

A．由题知导体棒匀速切割磁感线，根据几何关系切割长度为
$L=2x\tan\theta$，$x=vt$
则产生的感应电动势为
$E=2Bv^{2}t\tan\theta$
由题图可知电容器直接与电源相连，则电容器的电荷量为
$Q=CE=2BCv^{2}t\tan\theta$
则流过导体棒的电流
$I=\frac{\Delta Q}{\Delta t}=2BCv^{2}\tan\theta$
A 正确；

B．当金属棒到达 $x_{0}$ 处时，导体棒产生的感应电动势为
$E'=2Bvx_{0}\tan\theta$
则电容器的电荷量为
$Q=CE'=2BCvx_{0}\tan\theta$
B 错误；

D．由于导体棒做匀速运动则
$F=F_{\text{安}}=BIL$
由选项 A 可知流过导体棒的电流 $I$ 恒定，但 $L$ 与 $t$ 成正比，则 $F$ 为变力，再根据力做功的功率公式
$P=Fv$
可看出 $F$ 为变力，$v$ 不变则功率 $P$ 随力 $F$ 变化而变化；D 错误；

故选 A。
}

\item
%% number: 8
%% paperName: 2021年河北高考第 8 题 6 分
%% typeId: 2
%% type: 多选题
%% chapter: 交流电
%% point: 变压器
%% method: 无
%% score: 6
%% degree: 450
%% duplicateId: 0
%% body:
如图，发电机的矩形线圈长为 $2L$、宽为 $L$，匝数为 $N$，放置在磁感应强度大小为 $B$ 的匀强磁场中，理想变压器的原、副线圈匝数分别为 $n_{0}$、$n_{1}$ 和 $n_{2}$，两个副线圈分别接有电阻 $R_{1}$ 和 $R_{2}$，当发电机线圈以角速度 $\omega$ 匀速转动时，理想电流表读数为 $I$，不计线圈电阻，下列说法正确的是\xzanswer{BC}
\onepicture[w=8.0cm]{figs/08.png}
\fourchoices[answer=BC]
{通过电阻 $R_{2}$ 的电流为 $\frac{n_{1}I}{n_{2}}$}
{电阻 $R_{2}$ 两端的电压为 $\frac{n_{2}IR_{1}}{n_{1}}$}
{$n_{0}$ 与 $n_{1}$ 的比值为 $\frac{\sqrt{2}NBL^{2}\omega}{IR_{1}}$}
{发电机的功率为 $\frac{\sqrt{2}NBL^{2}\omega I(n_{1}+n_{2})}{n_{0}}$}

%% answer: BC
%% memo:
\memoanswer{
AB．由题知理想电流表读数为 $I$，则根据欧姆定律
$U_{1}=IR_{1}$
根据变压器电压与匝数的关系有
$\frac{n_{0}}{n_{1}}=\frac{U_{0}}{U_{1}}$，$\frac{n_{0}}{n_{2}}=\frac{U_{0}}{U_{2}}$
代入数据有
$U_{0}=\frac{n_{0}}{n_{1}}IR_{1}$，$U_{2}=\frac{n_{2}}{n_{1}}IR_{1}$
再由欧姆定律有
$U_{2}=I_{2}R_{2}$
可计算出
$I_{2}=\frac{n_{2}R_{1}}{n_{1}R_{2}}I$
综上可知，A 错误、B 正确；

C．由于矩形线圈产生的交变电流直接输入原线圈，则有
$E_{\max}=NB2L^{2}\omega$，$U_{0}=\frac{E_{\max}}{\sqrt{2}}=\sqrt{2}NBL^{2}\omega$
由选项 AB 知
$U_{0}=\frac{n_{0}}{n_{1}}IR_{1}$
则
$\frac{n_{0}}{n_{1}}=\frac{\sqrt{2}NBL^{2}\omega}{IR_{1}}$
C 正确；

D．由于变压器为理想变压器则有
$P_{0}=P_{1}+P_{2}=U_{1}I+U_{2}I_{2}=I^{2}R_{1}+U_{2}I_{2}$
代入选项 ABC 公式有
$P_{0}=\frac{\sqrt{2}NBL^{2}\omega I}{n_{0}}\left(\frac{n_{1}^{2}R_{2}+n_{2}^{2}R_{1}}{n_{1}R_{2}}\right)$
由于矩形线圈产生的交变电流直接输入原线圈，则发电机的功率为 $P_{0}$，D 错误。故选 BC。
}

\item
%% number: 9
%% paperName: 2021年河北高考第 9 题 6 分
%% typeId: 1
%% type: 单选题
%% chapter: 曲线运动
%% point: 竖直平面圆周运动
%% method: 无
%% score: 6
%% degree: 450
%% duplicateId: 0
%% body:
如图，矩形金属框 MNQP 竖直放置，其中 MN、PQ 足够长，且 PQ 杆光滑，一根轻弹簧一端固定在 M 点，另一端连接一个质量为 $m$ 的小球，小球穿过 PQ 杆，金属框绕 MN 轴分别以角速度 $\omega$ 和 $\omega'$ 匀速转动时，小球均相对 PQ 杆静止，若 $\omega'>\omega$，则与以 $\omega$ 匀速转动时相比，以 $\omega'$ 匀速转动时\xzanswer{BD}
\onepicture[w=5.5cm]{figs/09.png}
\fourchoices[answer=BD]
{小球的高度一定降低}
{弹簧弹力的大小一定不变}
{小球对杆压力的大小一定变大}
{小球所受合外力的大小一定变大}

%% answer: BD
%% memo:
\memoanswer{
对小球受力分析，设弹力为 $T$，弹簧与水平方向的夹角为 $\theta$，则对小球竖直方向
$T\sin\theta=mg$
而
$T=k\left(\frac{MP}{\cos\theta}-l_{0}\right)$
可知 $\theta$ 为定值，$T$ 不变，则当转速增大后，小球的高度不变，弹簧的弹力不变。则 A 错误，B 正确；水平方向当转速较小时，杆对小球的弹力 $F_{N}$ 背离转轴，则
$T\cos\theta-F_{N}=m\omega^{2}r$
即
$F_{N}=T\cos\theta-m\omega^{2}r$
当转速较大时，$F_{N}$ 指向转轴
$T\cos\theta+F_{N}'=m\omega^{2}r$
即
$F_{N}'=m\omega^{2}r-T\cos\theta$
则因 $\omega'>\omega$，根据牛顿第三定律可知，小球对杆的压力不一定变大。则 C 错误；
根据
$F_{\text{合}}=m\omega^{2}r$
可知，因角速度变大，则小球受合外力变大。则 D 正确。

故选 BD。
}

\item
%% number: 10
%% paperName: 2021年河北高考第 10 题 6 分
%% typeId: 2
%% type: 多选题
%% chapter: 电场
%% point: 电势能电势差电场力做功
%% method: 无
%% score: 6
%% degree: 350
%% duplicateId: 0
%% body:
如图，四个电荷量均为 $q$（$q>0$）的点电荷分别放置于菱形的四个顶点，其坐标分别为 $(4l,0)$、$(-4l,0)$、$(0,y_{0})$ 和 $(0,-y_{0})$，其中 $x$ 轴上的两个点电荷位置固定，$y$ 轴上的两个点电荷可沿 $y$ 轴对称移动（$y_{0}\neq0$），下列说法正确的是\xzanswer{ACD}
\onepicture[w=5.0cm]{figs/10.png}
\fourchoices[answer=ACD]
{除无穷远处之外，菱形外部电场强度处处不为零}
{当 $y_{0}$ 取某值时，可使得菱形内部只存在两个电场强度为零的点}
{当 $y_{0}=8l$ 时，将一带负电的试探电荷由点 $(4l,5l)$ 移至点 $(0,-3l)$，静电力做正功}
{当 $y_{0}=4l$ 时，将一带负电的试探电荷放置在点 $(l,l)$ 处，其所受到的静电力方向与 $x$ 轴正方向成 $45^{\circ}$ 倾斜向上}

%% answer: ACD
%% memo:
\memoanswer{
A．根据场强叠加原理可知，除无穷远处之外，菱形外部电场强度处处不为零，选项 A 正确；

B．因为在 $x$ 轴上的两个点电荷在 O 点的合场强为零，在 $y$ 轴上的两电荷，无论 $y_{0}$ 取什么值，因为关于原点对称，则在 O 点的合场强也为零，在横轴和纵轴上除原点外，出现合场强为零的点，根据对称性可知，一定是成对出现的，关于原点对称，所以算上原点，合场强为零的点是奇数个，不会是 2 个，选项 B 错误；

C．由几何关系可知，坐标为 $(4l,5l)$ 的 A 点在第一象限内所在的虚像的垂直平分线的上方；坐标为 $(0,-3l)$ 的 B 点在第三象限内所在的虚像的垂直平分线的上方，且到达虚线的距离相等，由电势叠加可知，B 点的电势高于 A 点，则带负电的试探电荷在 A 点的电势能较大，从 A 点到 B 点电势能减小，可知电场力做正功，选项 C 正确；
\onepicture[w=3.0cm]{figs/10_an1.png}

D．若 $y_{0}=4l$，则四个点构成正方形，由对称可知在点 $(l,l)$ 处的场强一定沿着过该点与原点连线的方向上；在 $y$ 轴正向和 $x$ 正向上的点电荷在 $(l,l)$ 处的合场强
$E_{1}=2\frac{kq}{(\sqrt{9l^{2}+l^{2}})^{2}}\cdot\frac{2\sqrt{2}l-\sqrt{2}l}{\sqrt{9l^{2}+l^{2}}}=\frac{kq}{5\sqrt{5}l^{2}}$
在 $y$ 轴负向和 $x$ 负向上的点电荷在 $(l,l)$ 处的合场强
$E_{2}=2\frac{kq}{(\sqrt{25l^{2}+l^{2}})^{2}}\cdot\frac{2\sqrt{2}l+\sqrt{2}l}{\sqrt{25l^{2}+l^{2}}}=\frac{kq}{\frac{13}{3}\sqrt{13}l^{2}}<E_{1}$
\onepicture[w=4.5cm]{figs/10_an2.png}
可知 $(l,l)$ 点的场强沿着 MN 方向且与 $x$ 轴成 $45^{\circ}$ 角的方向向下，将一带负电的试探电荷放置在点 $(l,l)$ 处，其所受到的静电力方向与 $x$ 轴正方向成 $45^{\circ}$ 倾斜向上，选项 D 正确。

故选 ACD。
}

\item
%% number: 11
%% paperName: 2021年河北高考第 11 题 6 分
%% typeId: 4
%% type: 实验题
%% chapter: 电路
%% point: 伏安法测电阻
%% method: 无
%% score: 6
%% degree: 650
%% duplicateId: 0
%% body:
某同学研究小灯泡的伏安特性，实验室提供的器材有；小灯泡（$6.3\UV$，$0.15\UA$），直流电源（$9\UV$），滑动变阻器，量程合适的电压表和电流表，开关和导线若干，设计的电路如图~\ref{2021河北11a}所示。
\begin{enumerate}
	\item
	根据图~\ref{2021河北11a}，完成图~\ref{2021河北11b}中的实物连线\_\_\_\_\_\_；
	\item
	按照图~\ref{2021河北11a}连线后，闭合开关，小灯泡闪亮一下后熄灭，观察发现灯丝被烧断，原因可能是\_\_\_\_\_\_（单项选择，填正确答案标号）；
	
	A．电流表短路
	
	B．滑动变阻器的滑片接触不良
	
	C．滑动变阻器滑片的初始位置在 b 端
	\item
	更换小灯泡后，该同学正确完成了实验操作，将实验数据描点作图，得到 $I-U$ 图像，其中一部分如图~\ref{2021河北11c}所示，根据图像计算出 P 点对应状态下小灯泡的电阻为\_\_\_\_\_\_$\UO$（保留三位有效数字）。
\end{enumerate}
\threepicture[labelA=2021河北11a, labelB=2021河北11b, labelC=2021河北11c, numA=1, numB=2, numC=3, showlabel=false, wA=5.5cm, wB=6.5cm, wC=6.5cm, align=right]
{figs/11a.png}
{figs/11b.png}
{figs/11c.png}

%% answer: 
\jdanswer{
\begin{enumerate}
	\item
	见详解图
	\item
	C
	\item
	27.0
\end{enumerate}
}
%% memo:
\memoanswer{
（1）电流表负极与滑动变阻器的右端的 b 位置连接，如图
\onepicture[w=8.0cm]{figs/11_an.png}

（2）开关闭合，小灯泡闪亮一下后灯丝烧断，说明通过小灯泡的电流过大。

A．电流表内阻非常小，短路几乎不影响通过小灯泡的电流，与灯丝烧断无关，A 错误；

B．滑动变阻器滑片接触不良，无电流通过小灯泡，B 错误；

C．滑动变阻器的滑片开始时置于 b 端，小灯泡部分分压达到最大，通过电流最大，可能会烧断小灯泡灯丝，C 正确；

故选 C。

（3）根据小灯泡的伏安特性曲线可知在 P 点时的电压和电流分别为
$U=2\UV$，$I=74\UmA$
根据欧姆定律可知小灯泡的电阻为
$R=\frac{U}{I}=\frac{2}{0.074}\UO=27.0\UO$
}

\item
%% number: 12
%% paperName: 2021年河北高考第 12 题 9 分
%% typeId: 4
%% type: 实验题
%% chapter: 功和能
%% point: 功能原理
%% method: 无
%% score: 9
%% degree: 450
%% duplicateId: 0
%% body:
某同学利用图~\ref{2021河北12a}中的实验装置探究机械能变化量与力做功的关系，所用器材有：一端带滑轮的长木板、轻细绳、$50\Ug$ 的钩码若干、光电门 2 个、数字计时器、带遮光条的滑块（质量为 $200\Ug$，其上可放钩码）、刻度尺，当地重力加速度为 $9.80\Umsq$，实验操作步骤如下：
\begin{steps}
	\item
	安装器材，调整两个光电门距离为 $50.00\Ucm$，轻细绳下端悬挂 4 个钩码，如图~\ref{2021河北12a}所示；
	\item
	接通电源，释放滑块，分别记录遮光条通过两个光电门的时间，并计算出滑块通过两个光电门的速度；
	\item
	保持最下端悬挂 4 个钩码不变，在滑块上依次增加一个钩码，记录滑块上所载钩码的质量，重复上述步骤；
	\item
	完成 5 次测量后，计算出每次实验中滑块及所载钩码的总质量 $M$、系统（包含滑块、滑块所载钩码和轻细绳悬挂钩码）总动能的增加量 $\Delta E_{k}$ 及系统总机械能的减少量 $\Delta E$，结果如下表所示：
\end{steps}

\begin{center}
\begin{tblr}{|c|c|c|c|c|c|}
\hline
$M/\Ukg$ & 0.200 & 0.250 & 0.300 & 0.350 & 0.400 \\
\hline
$\Delta E_{k}/\UJ$ & 0.582 & 0.490 & 0.392 & 0.294 & 0.195 \\
\hline
$\Delta E/\UJ$ & 0.393 & 0.490 & & 0.686 & 0.785 \\
\hline
\end{tblr}
\end{center}

回答下列问题：
\begin{enumerate}
	\item
	实验中轻细绳所悬挂钩码重力势能的减少量为\_\_\_\_\_\_$\UJ$（保留三位有效数字）；
	\item
	步骤④中的数据所缺数据为\_\_\_\_\_\_；
	\item
	若 $M$ 为横轴，$\Delta E$ 为纵轴，选择合适的标度，在图~\ref{2021河北12b}中绘出 $\Delta E-M$ 图像\_\_\_\_\_\_；
\end{enumerate}

若系统总机械能的减少量等于克服摩擦力做功，则物块与木板之间的摩擦因数为\tkanswer[2em]{0.40}（保留两位有效数字）。

\twopicture[labelA=2021河北12a, labelB=2021河北12b, numA=1, numB=2, showlabel=false, wA=7.5cm, wB=5.5cm, align=right]
{figs/12a.png}
{figs/12b.png}

%% answer: 
\jdanswer{
\begin{enumerate}
	\item
	0.980
	\item
	0.588
	\item
	见详解图。
\end{enumerate}
}
%% memo:
\memoanswer{
（1）四个钩码重力势能的减少量为
$\Delta E_{\text{p}}=4mgL=4\times0.05\times9.8\times0.5\UJ=0.980\UJ$

（2）对滑块和钩码构成的系统，由能量守恒定律可知
$4mgL-W_{\text{f}}=\frac{1}{2}(4m+M)v_{2}^{2}-\frac{1}{2}(4m+M)v_{1}^{2}$
其中系统减少的重力势能为
$\Delta E_{\text{p}}=4mgL$
系统增加的动能为
$\Delta E_{\text{k}}=\frac{1}{2}(4m+M)v_{2}^{2}-\frac{1}{2}(4m+M)v_{1}^{2}$
系统减少的机械能为 $\Delta E=W_{\text{f}}$，则代入数据可得表格中减少的机械能为
$\Delta E_{4}=0.98-0.392=0.588\UJ$

（3）根据表格数据描点得 $\Delta E-M$ 的图像为
\onepicture[w=8.0cm]{figs/12_an.png}
根据做功关系可知
$\Delta E=\mu MgL$
则 $\Delta E-M$ 图像的斜率为
$k=\mu gL=\frac{0.785-0.393}{0.4-0.2}=1.96$
解得动摩擦因数为
$\mu=0.40$（0.38～0.42）
}

\item
%% number: 13
%% paperName: 2021年河北高考第 13 题 11 分
%% typeId: 6
%% type: 计算题
%% chapter: 运动和力
%% point: 动量守恒定律
%% method: 无
%% score: 11
%% degree: 550
%% duplicateId: 0
%% body:
如图，一滑雪道由 AB 和 BC 两段滑道组成，其中 AB 段倾角为 $\theta$，BC 段水平，AB 段和 BC 段由一小段光滑圆弧连接，一个质量为 $2\Ukg$ 的背包在滑道顶端 A 处由静止滑下，若 $1\Us$ 后质量为 $48\Ukg$ 的滑雪者从顶端以 $1.5\Ums$ 的初速度、$3\Umsq$ 的加速度匀加速追赶，恰好在坡底光滑圆弧的水平处追上背包并立即将其拎起，背包与滑道的动摩擦因数为 $\mu=\frac{1}{12}$，重力加速度取 $g=10\Umsq$，$\sin\theta=\frac{7}{25}$，$\cos\theta=\frac{24}{25}$，忽略空气阻力及拎包过程中滑雪者与背包的重心变化，求：
\begin{enumerate}
	\item
	滑道 AB 段的长度；
	\item
	滑雪者拎起背包时这一瞬间的速度。
\end{enumerate}
\onepicture[w=7.0cm, align=right]{figs/13.png}

%% answer: 
\jdanswer{
\begin{enumerate}
	\item
	$L=9\Um$
	\item
	$v=7.44\Ums$
\end{enumerate}
}
%% memo:
\memoanswer{
（1）设斜面长度为 $L$，背包质量为 $m_{1}=2\Ukg$，在斜面上滑行的加速度为 $a_{1}$，由牛顿第二定律有
$m_{1}g\sin\theta-\mu m_{1}g\cos\theta=m_{1}a_{1}$
解得
$a_{1}=2\Umsq$
滑雪者质量为 $m_{2}=48\Ukg$，初速度为 $v_{0}=1.5\Ums$，加速度为 $a_{2}=3\Umsq$，在斜面上滑行时间为 $t$，落后时间 $t_{0}=1\Us$，则背包的滑行时间为 $t+t_{0}$，由运动学公式得
$L=\frac{1}{2}a_{1}(t+t_{0})^{2}$
$L=v_{0}t+\frac{1}{2}a_{2}t^{2}$
联立解得
$t=2\Us$ 或 $t=-1\Us$（舍去）
故可得
$L=9\Um$

（2）背包和滑雪者到达水平轨道时的速度为 $v_{1}$、$v_{2}$，有
$v_{1}=a_{1}(t+t_{0})=6\Ums$
$v_{2}=v_{0}+a_{2}t=7.5\Ums$
滑雪者拎起背包的过程，系统在光滑水平面上外力为零，动量守恒，设共同速度为 $v$，有
$m_{1}v_{1}+m_{2}v_{2}=(m_{1}+m_{2})v$
解得
$v=7.44\Ums$
}

\item
%% number: 14
%% paperName: 2021年河北高考第 14 题 61 分
%% typeId: 6
%% type: 计算题
%% chapter: 磁场
%% point: 带电粒子在磁场中的运动
%% method: 无
%% score: 61
%% degree: 350
%% duplicateId: 0
%% body:
如图，一对长平行栅极板水平放置，极板外存在方向垂直纸面向外、磁感应强度大小为 $B$ 的匀强磁场，极板与可调电源相连，正极板上 O 点处的粒子源垂直极板向上发射速度为 $v_{0}$、带正电的粒子束，单个粒子的质量为 $m$、电荷量为 $q$，一足够长的挡板 OM 与正极板成 $37^{\circ}$ 倾斜放置，用于吸收打在其上的粒子，C、P 是负极板上的两点，C 点位于 O 点的正上方，P 点处放置一粒子靶（忽略靶的大小），用于接收从上方打入的粒子，CP 长度为 $L_{0}$，忽略栅极的电场边缘效应、粒子间的相互作用及粒子所受重力。$\sin37^{\circ}=\frac{3}{5}$。
\begin{enumerate}
	\item
	若粒子经电场一次加速后正好打在 P 点处的粒子靶上，求可调电源电压 $U_{0}$ 的大小；
	\item
	调整电压的大小，使粒子不能打在挡板 OM 上，求电压的最小值 $U_{\min}$；
	\item
	若粒子靶在负极板上的位置 P 点左右可调，则负极板上存在 H、S 两点（$CH\leqslant CP<CS$，H、S 两点未在图中标出）、对于粒子靶在 HS 区域内的每一点，当电压从零开始连续缓慢增加时，粒子靶均只能接收到 $n$（$n\geqslant2$）种能量的粒子，求 CH 和 CS 的长度（假定在每个粒子的整个运动过程中电压恒定）。
\end{enumerate}
\onepicture[w=7.0cm, align=right]{figs/14.png}

%% answer: 
\jdanswer{
\begin{enumerate}
	\item
	$U_{0}=\frac{B^{2}qL_{0}^{2}}{8m}-\frac{mv_{0}^{2}}{2q}$
	\item
	$U_{\min}=\frac{7mv_{0}^{2}}{18q}$
	\item
	$CH=\frac{10mv_{0}}{3qB}$；$CS\to\infty$
\end{enumerate}
}
%% memo:
\memoanswer{
（1）从 O 点射出的粒子在板间被加速，则
$U_{0}q=\frac{1}{2}mv^{2}-\frac{1}{2}mv_{0}^{2}$
粒子在磁场中做圆周运动，则半径
$r=\frac{L_{0}}{2}$
由
$qvB=m\frac{v^{2}}{r}$
解得
$U_{0}=\frac{B^{2}qL_{0}^{2}}{8m}-\frac{mv_{0}^{2}}{2q}$

（2）当电压有最小值时，当粒子穿过下面的正极板后，圆轨道与挡板 OM 相切，此时粒子恰好不能打到挡板上，则
\onepicture[w=6.5cm]{figs/14_an1.png}
从 O 点射出的粒子在板间被加速，则
$U_{\min}q=\frac{1}{2}mv^{2}-\frac{1}{2}mv_{0}^{2}$
粒子在负极板上方的磁场中做圆周运动
$qvB=m\frac{v^{2}}{r_{\min}}$
粒子从负极板传到正极板时速度仍减小到 $v_{0}$，则
$qv_{0}B=m\frac{v_{0}^{2}}{r}$
由几何关系可知
$2r_{\min}=\frac{r'}{\sin37^{\circ}}+r'$
联立解得
$v=\frac{4v_{0}}{3}$
$U_{\min}=\frac{7mv_{0}^{2}}{18q}$

（3）设粒子第一次经过电场加速，在负极板上方磁场区域偏转的轨迹半径为 $r_{0}$，若粒子在电场加速电压小于 $U_{\min}$，粒子穿过磁场在正极板下方磁场运动时，会被 OM 板吸收。则第一次出现能吸收到两种能量的位置（即 H 点），为粒子通过极板电压 $U_{\min}=\frac{7mv_{0}^{2}}{18q}$ 时，粒子第二次从上方打到负极板的位置（轨迹如图中蓝色线条所示）。由（2）的计算可知
$r_{1}=\frac{4mv_{0}}{3qB}$
则
$CH=4r-2r'=\frac{10mv_{0}}{3qB}$
极板电压大于 $U_{\min}=\frac{7mv_{0}^{2}}{18q}$ 时，粒子均不会被 OM 吸收，可以经过正极板下方磁场偏转，回到负极板上方磁场中，偏转后打在负极板上。则 H 点右方的点的粒子靶都可以接受到 $n$（$n\geqslant2$）种能量的粒子。即 $CS\to\infty$。
\onepicture[w=6.5cm]{figs/14_an2.png}
}

\item
%% number: 15
%% paperName: 2021年河北高考第 15 题 4 分
%% typeId: 3
%% type: 填空题
%% chapter: 气体性质
%% point: 热力学第一定律
%% method: 无
%% score: 4
%% degree: 750
%% duplicateId: 0
%% body:
两个内壁光滑、完全相同的绝热汽缸 A、B，汽缸内用轻质绝热活塞封闭完全相同的理想气体，如图~\ref{2021河北15a}所示，现向活塞上表面缓慢倒入细沙，若 A 中细沙的质量大于 B 中细沙的质量，重新平衡后，汽缸 A 内气体的内能\tkanswer{大于}（填“大于”、“小于”或“等于”）汽缸 B 内气体的内能，图~\ref{2021河北15b}为重新平衡后 A、B 汽缸中气体分子速率分布图像，其中曲线\tkanswer{①}（填图像中曲线标号）表示汽缸 B 中气体分子的速率分布规律。
\twopicture[labelA=2021河北15a, labelB=2021河北15b, numA=1, numB=2, showlabel=false, wA=3.5cm, wB=7.5cm, align=right]
{figs/15a.png}
{figs/15b.png}

%% answer: 
%% memo:
\memoanswer{
（1）对活塞分析有
$p=\frac{mg}{S}$
因为 A 中细沙的质量大于 B 中细沙的质量，故稳定后有 $p_{A}>p_{B}$；所以在达到平衡过程中外界对气体做功有 $W_{A}>W_{B}$
则根据
$\Delta U=W+Q$
因为气缸和活塞都是绝热的，故有
$\Delta U_{A}>\Delta U_{B}$
即重新平衡后 A 气缸内的气体内能大于 B 气缸内的气体内能；

（2）由图中曲线可知曲线②中分子速率大的分子数占总分子数百分比较大，即曲线②的温度较高，所以由前面分析可知 B 气缸温度较低，故曲线①表示气缸 B 中气体分子的速率分布。
}

\item
%% number: 16
%% paperName: 2021年河北高考第 16 题 8 分
%% typeId: 6
%% type: 计算题
%% chapter: 气体性质
%% point: 查理定律
%% method: 无
%% score: 8
%% degree: 550
%% duplicateId: 0
%% body:
某双层玻璃保温杯夹层中有少量空气，温度为 $27\UdC$ 时，压强为 $3.0\times10^{3}\UPa$。
\begin{enumerate}
	\item
	当夹层中空气的温度升至 $37\UdC$，求此时夹层中空气的压强；
	\item
	当保温杯外层出现裂隙，静置足够长时间，求夹层中增加的空气质量与原有空气质量的比值，设环境温度为 $27\UdC$，大气压强为 $1.0\times10^{5}\UPa$。
\end{enumerate}

%% answer: 
\jdanswer{
\begin{enumerate}
	\item
	$p_{2}=3.1\times10^{3}\UPa$
	\item
	$\frac{97}{3}$
\end{enumerate}
}
%% memo:
\memoanswer{
（1）由题意可知夹层中的气体发生等容变化，根据理想气体状态方程可知
$\frac{p_{1}}{T_{1}}=\frac{p_{2}}{T_{2}}$
代入数据解得
$p_{2}=3.1\times10^{3}\UPa$

（2）当保温杯外层出现裂缝后，静置足够长时间，则夹层压强和大气压强相等，设夹层体积为 $V$，以静置后的所有气体为研究对象有
$p_{0}V=p_{1}V_{1}$
解得
$V_{1}=\frac{100}{3}V$
则增加空气的体积为
$\Delta V=V_{1}-V=\frac{97}{3}V$
所以增加的空气质量与原有空气质量之比为
$\frac{\Delta m}{m}=\frac{\Delta V}{V}=\frac{97}{3}$
}

\item
%% number: 17
%% paperName: 2021年河北高考第 17 题 4 分
%% typeId: 3
%% type: 填空题
%% chapter: 机械振动
%% point: 弹簧振子
%% method: 无
%% score: 4
%% degree: 750
%% duplicateId: 0
%% body:
如图，一弹簧振子沿 $x$ 轴做简谐运动，振子零时刻向右经过 A 点，$2\Us$ 后第一次到达 B 点，已知振子经过 A、B 两点时的速度大小相等，$2\Us$ 内经过的路程为 $0.4\Um$。该弹簧振子的周期为 \tkanswer{4} $\Us$，振幅为 \tkanswer{0.2} $\Um$。
\onepicture[h=2.5cm, align=right]{figs/17.png}

%% answer: 
%% memo:
\memoanswer{
（1）根据简谐运动对称性可知，振子零时刻向右经过 A 点，$2\Us$ 后第一次到达 B 点，已知振子经过 A、B 两点时的速度大小相等，则 A、B 两点关于平衡位置对称，而振动经过了半个周期的运动，则周期为
$T=2t=4\Us$

（2）从 A 到 B 经过了半个周期的振动，路程为 $s=0.4\Um$，而一个完整的周期路程为 $0.8\Um$，为 4 个振幅的路程，有
$4A=0.8\Um$
解得振幅为
$A=0.2\Um$
}

\item
%% number: 18
%% paperName: 2021年河北高考第 18 题 8 分
%% typeId: 6
%% type: 计算题
%% chapter: 光学
%% point: 光的折射
%% method: 无
%% score: 8
%% degree: 650
%% duplicateId: 0
%% body:
将两块半径均为 $R$、完全相同的透明半圆柱体 A、B 正对放置，圆心上下错开一定距离，如图所示，用一束单色光沿半径照射半圆柱体 A，设圆心处入射角为 $\theta$，当 $\theta=60^{\circ}$ 时，A 右侧恰好无光线射出；当 $\theta=30^{\circ}$ 时，有光线沿 B 的半径射出，射出位置与 A 的圆心相比下移 $h$，不考虑多次反射，求：
\begin{enumerate}
	\item
	半圆柱体对该单色光的折射率；
	\item
	两个半圆柱体之间的距离 $d$。
\end{enumerate}
\onepicture[w=5.0cm, align=right]{figs/18.png}

%% answer: 
\jdanswer{
\begin{enumerate}
	\item
	$n=\frac{2\sqrt{3}}{3}$
	\item
	$d=\sqrt{2}\left(h-\frac{R}{2}\right)$
\end{enumerate}
}
%% memo:
\memoanswer{
（i）光从半圆柱体 A 射入，满足从光密介质到光疏介质，当 $\theta=60^{\circ}$ 时发生全反射，有
$\sin\theta=\frac{1}{n}$
解得
$n=\frac{2\sqrt{3}}{3}$

（ii）当入射角 $\theta=30^{\circ}$，经两次折射从半圆柱体 B 的半径出射，设折射角为 $r$，光路如图
\onepicture[w=5.5cm]{figs/18_an.png}
由折射定律有
$\sin\theta\cdot n=\sin r$
有几何关系有
$\tan r=\frac{h-R\sin\theta}{d}$
联立解得
$d=\sqrt{2}\left(h-\frac{R}{2}\right)$
}

\end{enumerate}

\end{document}
